% META: {"id": "TNSI-BDD-CO-006", "chapitre": "TNSI-BASES-DE-DONNEES", "type_objet": "corrige", "exercice_ref": "TNSI-BDD-EX-006", "capacites": ["T-BDD-02"], "capacites_codes": ["C4"], "mode_creation": "ex_nihilo", "sources_inspiration": [], "status": "generated"}
% BEGIN-VERIFY
% import sqlite3
% conn = sqlite3.connect(":memory:")
% conn.execute("CREATE TABLE Adherent(id INTEGER PRIMARY KEY, nom TEXT, cotisation INTEGER)")
% conn.executemany("INSERT INTO Adherent VALUES (?,?,?)", [
%     (1, "Amine", 120), (2, "Lina", 95), (3, "Omar", 120), (4, "Sara", 150)])
% conn.commit()
% conn.execute("BEGIN")
% conn.execute("UPDATE Adherent SET cotisation = 130 WHERE id = 1")
% assert conn.execute("SELECT cotisation FROM Adherent WHERE id = 1").fetchone()[0] == 130
% conn.execute("ROLLBACK")
% assert conn.execute("SELECT cotisation FROM Adherent WHERE id = 1").fetchone()[0] == 120
% conn.execute("BEGIN")
% conn.execute("UPDATE Adherent SET cotisation = 130 WHERE id = 1")
% conn.execute("COMMIT")
% assert conn.execute("SELECT cotisation FROM Adherent WHERE id = 1").fetchone()[0] == 130
% conn.execute("UPDATE Adherent SET cotisation = 120 WHERE id = 1")
% conn.commit()
% avant = conn.execute("SELECT nom FROM Adherent WHERE cotisation = 120").fetchall()
% conn.execute("CREATE INDEX idx_cotisation ON Adherent(cotisation)")
% assert conn.execute("SELECT nom FROM Adherent WHERE cotisation = 120").fetchall() == avant
% assert avant == [("Amine",), ("Omar",)]
% plan = str(conn.execute(
%     "EXPLAIN QUERY PLAN SELECT nom FROM Adherent WHERE cotisation = 120").fetchall())
% assert "idx_cotisation" in plan
% conn.execute("CREATE VIEW Annuaire AS SELECT id, nom FROM Adherent")
% assert [c[1] for c in conn.execute("PRAGMA table_info(Annuaire)")] == ["id", "nom"]
% assert conn.execute("SELECT COUNT(*) FROM Annuaire").fetchone()[0] == 4
% refuse = False
% try:
%     conn.execute("SELECT cotisation FROM Annuaire")
% except sqlite3.OperationalError:
%     refuse = True
% assert refuse
% END-VERIFY
\begin{corrige}{TNSI-BDD-EX-006}
\textbf{1.} L'association se fait sans ambiguïté :

\begin{center}
\begin{tabular}{ll}
\hline
mésaventure & service du SGBD \\
\hline
(a) coupure de courant & persistance \\
(b) modifications simultanées & gestion de la concurrence \\
(c) recherche lente & efficacité du traitement des données \\
(d) stagiaire voyant les cotisations & sécurisation des accès \\
\hline
\end{tabular}
\end{center}

\textbf{2.} \lstinline{BEGIN} ouvre une \emph{transaction} : les modifications
qui suivent sont provisoires et invisibles des autres utilisateurs.
\lstinline{COMMIT} les valide toutes d'un coup ; \lstinline{ROLLBACK} les annule
toutes. Le SGBD garantit qu'une transaction est indivisible : ou bien tout est
appliqué, ou bien rien.
\begin{sql}
BEGIN;
UPDATE Adherent SET cotisation = 130 WHERE id = 1;
ROLLBACK;
\end{sql}
Après le \lstinline{ROLLBACK}, la cotisation d'Amine vaut de nouveau $120$. Entre
le \lstinline{BEGIN} et le \lstinline{ROLLBACK}, la connexion qui a fait la
modification lit bien $130$ — mais elle est seule à la voir. C'est exactement ce
qui manquait au tableur : la seconde secrétaire y écrasait un travail dont rien
ne lui signalait l'existence.

\textbf{3.} \lstinline{CREATE INDEX idx_cotisation ON Adherent(cotisation);}

L'index ne change \textbf{pas} le résultat : la requête renvoie
\lstinline{Amine} et \lstinline{Omar} avant comme après. Il change le
\emph{chemin} suivi pour l'obtenir. Sans index, le SGBD parcourt toute la table ;
avec index, il consulte une structure ordonnée et va directement aux lignes
concernées. On peut le vérifier avec \lstinline{EXPLAIN QUERY PLAN}, qui
mentionne alors \lstinline{idx_cotisation}.

C'est un point de méthode important : optimiser, ici, ne consiste pas à réécrire
la requête, mais à donner au SGBD un moyen supplémentaire de la satisfaire. Le
programme qui interroge la base n'a pas une ligne à changer.

\textbf{4.} \lstinline{CREATE VIEW Annuaire AS SELECT id, nom FROM Adherent;}

La vue contient les $4$ lignes de la table — elle ne filtre pas les lignes, mais
les colonnes. Si le stagiaire demande \lstinline{SELECT cotisation FROM Annuaire},
le SGBD refuse : la colonne n'existe pas dans la vue. Il ne s'agit donc pas d'un
affichage masqué, mais d'une donnée hors de portée.

\textbf{5.} Deux exemples suffisent.

Pour (a), un fichier est écrit par le système d'exploitation : si la machine
s'arrête au milieu d'une écriture, le fichier reste dans un état intermédiaire, à
moitié ancien et à moitié nouveau, et rien ne permet de savoir lequel. Un SGBD
tient un \emph{journal} des transactions avant de toucher aux données : au
redémarrage, il rejoue ce qui était validé et annule le reste. Cette garantie ne
peut pas être ajoutée par-dessus un fichier ordinaire, elle doit être présente au
moment de l'écriture.

Pour (d), un fichier partagé se donne en entier ou pas du tout : les droits du
système portent sur le fichier, pas sur une colonne. Restreindre la vue du
stagiaire supposerait de maintenir une seconde copie sans les cotisations — donc
deux vérités à garder synchronisées, ce qui est précisément le problème que la
base était censée résoudre.
\end{corrige}
